%------------------------------------------------------------------------------%

\usepackage{integracao_e_series}
\usepackage{ulem}

\title{Integração por Substituição Trigonométrica}
\subtitle{2 -- Substituição pela tangente}

%------------------------------------------------------------------------------%
\begin{document}

\addtitle

%------------------------------------------------------------------------------%
\section{Integração por Substituição Trigonométrica}

%------------------------------------------------------------------------------%
\begin{frame}{Substituições Trigonométricas}
\addtext{16}{0.35,2.2}{
\renewcommand*{\arraystretch}{2.3}
\begin{tabular}{c<{\;}>{\;}c<{\;}>{\;}c<{\;}>{\;}c} \hline
  Expressão
  & Substituição
  & Intervalo
  & Identidade \\ \hline
  \(\sqrt{a^2-x^2}\)
  & \(x=a\sin(\theta)\)
  & \(-\frac{\pi}{2} \leq \theta \leq \frac{\pi}{2}\)
  & \(1-\sin^2(\theta) = \cos^2(\theta)\) \\
  {\color{blue}\(\sqrt{a^2+x^2}\)}
  & {\color{blue}\(x=a\tan(\theta)\)}
  & \(-\frac{\pi}{2}< \theta < \frac{\pi}{2} \)
  & \(1+\tan^2(\theta)= \sec^2(\theta)\) \\
  \(\sqrt{x^2-a^2}\)
  & \(x=a\sec(\theta)\)
  & \(0 \leq \theta < \frac{\pi}{2} \) ou \(\pi \leq \theta < \frac{3\pi}{2} \)
  & \(\sec^2(\theta)-1= \tan^2(\theta)\) \\ \hline
\end{tabular}}
\end{frame}

% 1
%------------------------------------------------------------------------------%
%------------------------------------------------------------------------------%
\addtocounter{exemplo}{1}
\section{Exemplo \theexemplo}

%------------------------------------------------------------------------------%
\begin{frame}{Exemplo \theexemplo}
  Encontre \(\int \frac{dx}{x^2\sqrt{x^2+4}}\)

  \pause
  \vspace{0.5\baselineskip}
  Note
  \(
    \sqrt{x^2+4} = \sqrt{2^2+x^2}
  \)

  \pause
  \vspace{\baselineskip}
  Substituição inversa
  \[
    x = 2\tan(\theta)
    \qquad\qquad
    dx = 2\sec^2(\theta) d\theta
    \qquad\text{com}\qquad
    -\frac{\pi}{2} < \theta < \frac{\pi}{2}
  \]
\end{frame}

%------------------------------------------------------------------------------%
\begin{frame}{Exemplo \theexemplo}
  \begin{align*}
    \onslide<+->{F & = \int \frac{1}{x^2\sqrt{x^2+4}}dx
                   & & x = 2\tan(\theta)                       \\[2mm]}
    \onslide<+->{  & = \int \frac{1}{4\tan^2(\theta) \sqrt{4\tan^2(\theta)+4}}\,2\sec^2(\theta) d\theta      \\[2mm]}
    \onslide<+->{  & = \int \frac{\sec^2(\theta)}{2\tan^2(\theta) \sqrt{4\big(1+ \tan^2(\theta)\big)}} d\theta \\[2mm]}
    \onslide<+->{  & = \int \frac{\sec^2(\theta)}{4\tan^2(\theta) \sqrt{1+ \tan^2(\theta)}} d\theta
                   & & 1 + \tan^2(\theta) = \sec^2(\theta) \\[2mm]}
    \onslide<+->{  & = \frac{1}{4} \int \frac{\sec^2(\theta)}{\tan^2(\theta) \sqrt{\sec^2(\theta)}} d\theta}
    %
  \end{align*}
\end{frame}

%------------------------------------------------------------------------------%
\begin{frame}{Exemplo \theexemplo}
  \begin{align*}
    \onslide<+->{F & = \frac{1}{4} \int \frac{\sec^2(\theta)}{\tan^2(\theta) \sqrt{\sec^2(\theta)}} d\theta\\[3mm]}
%    & = \int \frac{\sec^2(\theta)}{\tan^2(\theta) \sec(\theta)}  d\theta \\[2mm]
    \onslide<+->{  & = \frac{1}{4} \int \frac{\sec(\theta)}{\tan^2(\theta)} d\theta              \\[3mm]}
    \onslide<+->{  & = \frac{1}{4} \int \frac{1}{\cos(\theta)}
                                   \left(\frac{\cos(\theta)}{\sin(\theta)}\right)^{\!\!2} d\theta \\[3mm]}
    \onslide<+->{  & = \frac{1}{4} \int \frac{\cos(\theta)}{\sin^2(\theta)} d\theta}
  \end{align*}
\end{frame}

%------------------------------------------------------------------------------%
\begin{frame}{Exemplo \theexemplo}
  Integrar
  \(
    F = \frac{1}{4} \int \frac{\cos(\theta)}{\sin^2(\theta)} d\theta
  \)

  \pause
  \vspace{0.5\baselineskip}
  Substituição simples
  \[
    u = \sin(\theta)
    \qquad
    du = \cos(\theta) d\theta
  \]

  \pause
  Temos
  \[
    F
    = \frac{1}{4} \int \frac{du}{u^2}
    \pause
    = \frac{1}{4} \int u^{-2} du
    \pause
    = \frac{1}{4} \left(\frac{u^{-1}}{-1} + C\right)
    \pause
    = \frac{-1}{4\sin(\theta)} + C_1
  \]
\end{frame}

%------------------------------------------------------------------------------%
\begin{frame}{Exemplo \theexemplo}
  \[
    x= 2 \tan(\theta) \qquad\Rightarrow\qquad \tan(\theta) = \frac{x}{2}
  \]

  \addgraphic{6}{10,3}{exemplo_triangulo_2}

  \pause
  \[
    \frac{1}{\sin(\theta)}
    = \csc(\theta)
    \pause
    = \frac{\text{hipotenusa}}{\text{cateto oposto}}
    \pause
    = \frac{\sqrt{x^2+4}}{x}
  \]

  \pause
  \begin{align*}
    \onslide<+->{F & = \int \frac{dx}{x^2\sqrt{x^2+4}}  \\[2mm]}
    \onslide<+->{  & = \frac{-1}{4\sin(\theta)} + C \\[2mm]}
    \onslide<+->{  & = -\frac{\sqrt{x^2+4}}{4x} + C}
  \end{align*}
\end{frame}

% 2
%------------------------------------------------------------------------------%
%------------------------------------------------------------------------------%
\addtocounter{exemplo}{1}
\section{Exemplo \theexemplo}

%------------------------------------------------------------------------------%
\begin{frame}{Exemplo \theexemplo}
  \onslide<1->{Encontre \(\int \frac{x}{\sqrt{x^2+4}} dx\)}

  \vspace{0.5\baselineskip}
  \only<1|handout:0>{\vphantom{Substituição trigonométrica \(x = 2 \tan(\theta)\)}}
  \only<2|handout:0>{Substituição trigonométrica \(x = 2 \tan(\theta)\)}
  \only<3:->{\sout{Substituição trigonométrica \(x = 2 \tan(\theta)\)}}

  \vspace{0.5\baselineskip}
  \onslide<4->{Substituição simples \(u=x^2+4 \qquad du=2x dx\)}
  \[
    \onslide<5->{\int \frac{x}{\sqrt{x^2+4}} dx
                 = \frac{1}{2} \int \frac{du}{\sqrt{u}}}
    \onslide<6->{= \sqrt{u}+ C}
    \onslide<7->{= \sqrt{x^2+4}+C}
  \]
\end{frame}

% 3
%------------------------------------------------------------------------------%
%------------------------------------------------------------------------------%
\addtocounter{exemplo}{1}
\section{Exemplo \theexemplo}

%------------------------------------------------------------------------------%
\begin{frame}{Exemplo \theexemplo}
  Calcule \(\int \frac{dt}{\sqrt{t^2-6t+13}}\)

  \pause
  \vspace{\baselineskip}
  Completando quadrado temos
  \[
    t^2 - 6t + 13
    \pause
    \;=\; (t-3)^2 - 9 + 13
    \pause
    \;=\; (t-3)^2 + 4
  \]

  \pause
  \[
    F = \int \frac{dt}{\sqrt{t^2-6t+13}}
    \pause
      = \int \frac{dt}{\sqrt{(t-3)^2 + 4}}
  \]
\end{frame}

%------------------------------------------------------------------------------%
\begin{frame}{Exemplo \theexemplo\ -- Substituição Simples}
  \onslide<+->{\[
    u = t - 3 \qquad\qquad du = dt
  \]}
  \vspace{-\baselineskip}
  \begin{align*}
    \onslide<+->{F & = \int \frac{dt}{\sqrt{(t-3)^2+4}}   \\[2mm]}
    \onslide<+->{  & = \int \frac{du}{\sqrt{u^2+4}}       \\[2mm]}
    \onslide<+->{  & = \int \frac{du}{\sqrt{u^2+2^2}}}
  \end{align*}
\end{frame}

%------------------------------------------------------------------------------%
\begin{frame}{Exemplo \theexemplo\ -- Substituição Trigonométrica}
  \onslide<+->{\[
    u  = 2\tan(\theta) \quad\Rightarrow\quad
    du = 2\sec^2(\theta) d\theta
  \]}
  \begin{align*}
    \onslide<+->{F & = \int \frac{du}{\sqrt{u^2+2^2}}                                  \\[2mm]}
    \onslide<+->{  & = \int \frac{2\sec^2(\theta) }{\sqrt{4\tan^2(\theta)+4}}d\theta}
  \end{align*}
\end{frame}

%------------------------------------------------------------------------------%
\begin{frame}{Exemplo \theexemplo\ -- Substituição Trigonométrica}
  \begin{align*}
    \onslide<+->{F & = \int \frac{2\sec^2(\theta) }{\sqrt{4\tan^2(\theta)+4}}d\theta   \\[2mm]}
    \onslide<+->{  & = \int \frac{2\sec^2(\theta) }{2\sqrt{\tan^2(\theta)+1}}d\theta
      & \tan^2(\theta)+1= \sec^2(\theta)                                  \\[2mm]}
    \onslide<+->{  & = \int \frac{\sec^2(\theta)}{\sec(\theta)}d\theta                 \\[2mm]}
    \onslide<+->{  & = \int \sec(\theta)d\theta                                        \\[2mm]}
    \onslide<+->{  & = \ln \abs{\sec(\theta)+ \tan(\theta)} + K}
  \end{align*}
\end{frame}

%------------------------------------------------------------------------------%
\begin{frame}{Exemplo \theexemplo\ -- Desfazendo as Transformações}
  \onslide<+->{\[
    u = 2\tan(\theta)
    \quad\Rightarrow\quad
    \tan(\theta) = \frac{u}{2} = \frac{t-3}{2}
  \]}

  \addgraphic{5.5}{10,3}{exemplo_triangulo_5}

  \onslide<+->{Hipotenusa}
  \begin{align*}
    \onslide<2->{\sqrt{(t-3)^2+2^2}
                 & = \sqrt{t^2-6t+9+4} \\}
    \onslide<+->{& = \sqrt{t^2-6t+13}}
  \end{align*}
  \onslide<+->{
  \[
    \sec(\theta)=\frac{\sqrt{t^2-6t+13}}{2}
  \]}
\end{frame}

%------------------------------------------------------------------------------%
\begin{frame}{Exemplo \theexemplo\ -- Desfazendo as Transformações}
  \begin{align*}
    \onslide<+->{F & = \int \frac{dt}{\sqrt{t^2-6t+13}}         \\[4mm]}
    \onslide<+->{  & = \ln \abs{\sec(\theta)+ \tan(\theta)} + K \\[4mm]}
    \onslide<+->{  & = \ln \abs{\frac{\sqrt{t^2-6t+13}}{2} + \frac{t-3}{2}} + K \\[4mm]}
    \onslide<+->{  & = \ln \abs{\sqrt{t^2-6t+13} + t-3} - \ln(2) + K \\[4mm]}
    \onslide<+->{  & = \ln \abs{\sqrt{t^2-6t+13} + t-3} + K_1}
  \end{align*}
\end{frame}

%------------------------------------------------------------------------------%
\section{Lista Mínima}

%------------------------------------------------------------------------------%
\begin{frame}{Lista Mínima}
  Estudar a Seção 4.4 da Apostila

  Exercícios:

  \vfill
  Atenção: A prova é baseada no livro, não nas apresentações
\end{frame}

%------------------------------------------------------------------------------%
\end{document}
%------------------------------------------------------------------------------%
